\chapter{Curves}\label{ch:curves}

Curves are the starting point of almost every model: they are the profiles that are revolved,
extruded and swept into surfaces, the sections of lofts, the boundaries of Coons patches and the
trim loops of faces. In \NF{} every curve -- a straight line as much as a fitted free-form curve
-- is an instance of one class, \apiref{NurbsCurve}: a non-uniform rational B-spline in
three-dimensional space. This chapter first presents the functions that create the common curve
types exactly, then the \code{NurbsCurve} class with all its members, and finally the operations
that build new curves from existing ones (joining, extension, offsets, fillets) and the analysis
functions (curvature, arc length, closest points).

All curve functions return \emph{new} objects; a curve is never modified in place
(\cref{ch:concepts}). Angles are in radians throughout the Python and C++ API.

\section{Curve primitives}\label{sec:curve-primitives}

\Cref{tab:curve-primitives} lists the functions that construct curves from a few geometric
parameters. Lines, polylines, polygons, rectangles, arcs and circles are \emph{exact}: the NURBS
represents the geometric object up to floating-point rounding. Helices and spirals are not
representable as NURBS and are approximated to a tolerance.

\begin{table}[htbp]
\centering\small
\begin{tabularx}{\linewidth}{@{}l>{\raggedright\arraybackslash}Xll@{}}
\toprule
Function & Result & Degree & Rational \\
\midrule
\apiref{make_line}\code{(p0, p1)} & segment from \code{p0} to \code{p1} & 1 & no \\
\apiref{make_polyline}\code{(points)} & open polyline through the points, chord-length knots & 1 & no \\
\apiref{make_polygon}\code{(center, normal, radius, n, x_axis)} & closed regular $n$-gon inscribed in a circle & 1 & no \\
\apiref{make_rectangle_curve}\code{(origin, edge_u, edge_v)} & closed parallelogram outline & 1 & no \\
\apiref{make_arc}\code{(center, x_axis, y_axis, radius, start, end)} & circular arc between two angles & 2 & yes \\
\apiref{make_arc_3pt}\code{(p0, p1, p2)} & arc from \code{p0} through \code{p1} to \code{p2} & 2 & yes \\
\apiref{make_circle}\code{(center, normal, radius)} & full circle & 2 & yes \\
\apiref{make_circle_3pt}\code{(p0, p1, p2)} & full circle through three points & 2 & yes \\
\apiref{make_helix}\code{(axis_point, axis, start, pitch, turns, ...)} & cylindrical or conical helix, planar spiral & 5 & no \\
\bottomrule
\end{tabularx}
\caption{Curve primitives. All points and vectors accept \code{Vec3} objects, tuples or lists.}
\label{tab:curve-primitives}
\end{table}

The following script creates one curve of each kind and prints its degree, number of control
points, whether it is rational, its parameter domain and its length:

\example{05_curve_primitives.py}

Output:

\begin{lstlisting}[style=text]
line         degree=1 n_ctrl= 2 rational=False domain=(0.0, 1.0) length=10.000000
polyline     degree=1 n_ctrl= 4 rational=False domain=(0.0, 1.0) length=11.000000
arc 120 deg  degree=2 n_ctrl= 5 rational=True  domain=(0.0, 1.0) length=10.471976
arc 3 pt     degree=2 n_ctrl= 5 rational=True  domain=(0.0, 1.0) length=15.707963
circle       degree=2 n_ctrl= 9 rational=True  domain=(0.0, 1.0) length=31.415927
circle 3 pt  degree=2 n_ctrl= 9 rational=True  domain=(0.0, 1.0) length=6.283185
hexagon      degree=1 n_ctrl= 7 rational=False domain=(0.0, 1.0) length=12.000000
rectangle    degree=1 n_ctrl= 5 rational=False domain=(0.0, 1.0) length=12.000000
circle radius error: 1.8e-15
helix: domain=(0.0, 3.0), end point=Vec3(1.000000, -0.000000, 6.000000), length=19.781449
exact helix length 3*sqrt((2 pi)^2 + 2^2) = 19.781450
conical helix end=Vec3(0.250000, -0.000000, 6.000000), spiral end=Vec3(3.000000, -0.000000, 0.000000)
\end{lstlisting}

\subsection{Lines, polylines, polygons and rectangles}

\code{make_line(p0, p1)} returns the degree-1 B-spline with the two control points \code{p0} and
\code{p1} on the domain $[0,1]$; the parameter is proportional to the distance from \code{p0}.

\code{make_polyline(points)} connects two or more points by straight segments. The knots are
chosen by \emph{chord length}: the polyline through $(0,0,0)$, $(4,0,0)$, $(4,3,0)$, $(0,3,0)$
(segment lengths 4, 3 and 4) has the knot vector $\{0, 0, \tfrac{4}{11}, \tfrac{7}{11}, 1,
1\}$, so the parameter is proportional to arc length along the whole polyline. The interior
vertices are corners (the curve is only $C^0$ there); they can be rounded with
\apiref{fillet_corners} (\cref{sec:curve-fillets}).

\code{make_polygon(center, normal, radius, n, x_axis)} creates a closed regular polygon with
$n$ vertices on the circle of the given radius in the plane through \code{center} perpendicular
to \code{normal}. The first vertex lies in the direction \code{x_axis} (projected into the
plane). If \code{x_axis} is omitted (the default is the zero vector), an arbitrary but
deterministic perpendicular direction is used -- for the normal $(0,0,1)$ this is the $+y$
axis, so the first vertex of a hexagon of radius 2 is $(0,2,0)$. Pass \code{x_axis} whenever
the orientation of the polygon matters. The polygon has $n+1$ control points; the last equals
the first.

\code{make_rectangle_curve(origin, edge_u, edge_v)} is the closed outline through the four
corners $\mathbf o$, $\mathbf o+\mathbf e_u$, $\mathbf o+\mathbf e_u+\mathbf e_v$,
$\mathbf o+\mathbf e_v$ and back to $\mathbf o$ (with $\mathbf o$ = \code{origin}, $\mathbf
e_u$ = \code{edge_u}, $\mathbf e_v$ = \code{edge_v}), with chord-length knots. The two edge
vectors need not be perpendicular, so any parallelogram can be created. Do not confuse it with
\apiref{make_rectangle}, which returns the planar \emph{surface} patch (\cref{ch:surfaces}).

\subsection{Arcs and circles}\label{sec:arcs}

Circular arcs are represented exactly as rational quadratic NURBS (Piegl \& Tiller,
algorithm~A7.1). An arc is split into segments of at most $90^\circ$; each segment has three
control points, the middle one with weight $\cos(\Delta\theta/2)$. A full circle therefore has
four segments, nine control points and the knot vector
$\{0,0,0,\tfrac14,\tfrac14,\tfrac12,\tfrac12,\tfrac34,\tfrac34,1,1,1\}$ with the weights
$1, \tfrac{\sqrt2}{2}, 1, \tfrac{\sqrt2}{2}, \dots$ The script above confirms the exactness: the
largest deviation of the circle from radius 5 is $1.8\cdot10^{-15}$.

\begin{description}[style=nextline]
\item[\code{make_arc(center, x_axis, y_axis, radius, start_angle, end_angle)}]
  The arc of the circle with the given centre and radius in the plane spanned by \code{x_axis}
  and \code{y_axis}, running from \code{start_angle} to \code{end_angle}; angles are measured
  from \code{x_axis} towards \code{y_axis}. The point at angle $\theta$ is
  $\text{center} + r(\cos\theta\,\hat x + \sin\theta\,\hat y)$. The axes are normalised; they
  should be perpendicular. With \code{start = 0} and \code{end = math.radians(120)} the arc ends
  at $(-2.5, 4.330127, 0)$ for radius 5.
\item[\code{make_arc_3pt(p0, p1, p2)}]
  The arc from \code{p0} through \code{p1} to \code{p2}. The three points must not be collinear.
\item[\code{make_circle(center, normal, radius)}]
  The full circle, oriented counter-clockwise about \code{normal} (right-hand rule). The start
  point lies on the same deterministic perpendicular of the normal as for \code{make_polygon}: for
  the normal $(0,0,1)$ the circle starts at $\text{center} + r\,(0,1,0)$ and reaches
  $\text{center} + r\,(-1,0,0)$ at $u = 0.25$. If the start point matters (for example for a
  sweep or a loft whose sections must line up), use \code{make_arc} with explicit axes and the
  angles $0$ and $2\pi$, or \code{make_circle_3pt}.
\item[\code{make_circle_3pt(p0, p1, p2)}]
  The full circle through three points; it starts at \code{p0} and runs through \code{p1} and
  \code{p2} in this order, which fixes both start point and orientation.
\end{description}

\begin{note}[Parameter and arc length on a circle]
The quarter points of a circle lie exactly at $u = 0, \tfrac14, \tfrac12, \tfrac34$. Within a
quarter, however, the rational parametrisation is \emph{not} proportional to arc length (no
rational parametrisation of the circle is). Use \apiref{parameter_at_length}
(\cref{sec:curve-analysis}) to find points at equal distances along the circle.
\end{note}

\subsection{Helices and spirals}\label{sec:helix}

\begin{lstlisting}[style=text]
make_helix(axis_point, axis, start, pitch, turns, end_radius=-1.0,
           left_handed=False, tolerance=1e-7) -> NurbsCurve
\end{lstlisting}

\apiref{make_helix} creates a helix around the axis through \code{axis_point} with direction
\code{axis}. The curve starts at the point \code{start}; its distance from the axis is the start
radius, and the foot of \code{start} on the axis is the height zero. Each turn advances by
\code{pitch} along the axis. The parameter of the curve is the number of turns: the domain is
$[0, \text{turns}]$, and $u = 1$ is the end of the first turn.

\begin{itemize}
\item \code{end_radius < 0} (default): cylindrical helix with constant radius.
\item \code{end_radius >= 0}: the radius changes linearly from the start radius to
  \code{end_radius} -- a \emph{conical helix}. In the example the conical helix ends at radius
  0.25.
\item \code{pitch = 0} together with \code{end_radius}: a planar \emph{Archimedean spiral} (the
  radius grows linearly with the angle).
\item \code{left_handed=True} turns clockwise about \code{axis} instead of counter-clockwise; at
  a quarter turn the right-handed helix of the example is at $(0, 1, 0.5)$, the left-handed one
  at $(0, -1, 0.5)$.
\end{itemize}

The helix is approximated by a $C^2$ quintic B-spline whose distance from the exact curve is at
most \code{tolerance}. The default tolerance of $10^{-7}$ is small enough for nearly all
purposes; the length of the three-turn helix in the example agrees with the exact value
$3\sqrt{(2\pi)^2+2^2}$ to $10^{-6}$. Use a smaller tolerance for very large helices or when the
curvature and torsion (\cref{sec:curve-analysis}) must be accurate to many digits.

\section{The NurbsCurve class}\label{sec:nurbscurve}

\subsection{Definition}

A NURBS curve of degree $p$ with control points $\mathbf P_0,\dots,\mathbf P_n$, weights
$w_i > 0$ and a knot vector $U = \{u_0, \dots, u_{n+p+1}\}$ is
\[
  \mathbf C(u) = \frac{\sum_{i=0}^{n} N_{i,p}(u)\, w_i\, \mathbf P_i}
                      {\sum_{i=0}^{n} N_{i,p}(u)\, w_i},
  \qquad u \in [u_p, u_{n+1}],
\]
where $N_{i,p}$ are the B-spline basis functions. If all weights are equal the curve is a
(non-rational) B-spline. The knot vector must be non-decreasing, contain $n + p + 2$ values and
no interior knot may have a multiplicity larger than $p$. At a knot of multiplicity $k$ the curve
is $C^{p-k}$ continuous. The background is explained in \cref{ch:concepts}.

\subsection{Construction}

\begin{lstlisting}[style=text]
NurbsCurve(degree, knots, control_points, weights=[])
NurbsCurve.from_points(degree, control_points)
\end{lstlisting}

The constructor takes the degree, the full knot vector (a list of floats), the control points as
an array of shape $(n+1, 3)$ (a NumPy array or a nested list) and optionally the weights. An empty
weight list means a non-rational curve. The constructor validates everything and raises
\code{GeometryError} with a precise message if, for example, the knot vector has the wrong size,
decreases somewhere or a weight is not positive.

\apiref{NurbsCurve.from_points} is a shortcut for a non-rational B-spline with a \emph{clamped
uniform} knot vector on $[0,1]$: the first and last $p+1$ knots are 0 and 1, the interior knots
are equally spaced. For six control points and degree 3 the knot vector is
$\{0,0,0,0,\tfrac13,\tfrac23,1,1,1,1\}$. A clamped curve starts at the first and ends at the last
control point.

\subsection{Properties}

\begin{tabularx}{\linewidth}{@{}l>{\raggedright\arraybackslash}X@{}}
\toprule
Property & Meaning \\
\midrule
\code{degree} & polynomial degree $p$ \\
\code{knots} & the full knot vector as a list \\
\code{control_points} & NumPy array of shape $(n+1, 3)$ (a copy) \\
\code{weights} & list of $n+1$ weights (all 1 for a non-rational curve) \\
\code{domain} & tuple $(u_p, u_{n+1})$, the valid parameter range \\
\code{is_rational} & \code{True} if the weights are not all equal \\
\code{is_clamped} & \code{True} if the first and last knot have multiplicity $p+1$ \\
\bottomrule
\end{tabularx}

\subsection{Evaluation and derivatives}

\begin{description}[style=nextline]
\item[\code{c.evaluate(u)} or \code{c(u)}] The point $\mathbf C(u)$ as a \code{Vec3}. Evaluation
  is done in homogeneous coordinates (Piegl \& Tiller A4.1), so rational curves are exact.
\item[\code{c.derivatives(u, order=1)}] A list of \code{order + 1} vectors
  $[\mathbf C(u), \mathbf C'(u), \dots, \mathbf C^{(k)}(u)]$ -- derivatives of any order with
  respect to the parameter, including the quotient rule for rational curves (A4.2).
\item[\code{c.tangent(u)}] The unit tangent $\mathbf C'(u)/|\mathbf C'(u)|$.
\item[\code{c.sample(n=100)}] A NumPy array of shape $(n, 3)$ with the points at $n$ equally
  spaced parameters including both ends -- convenient for plotting with Matplotlib.
\end{description}

Parameters outside the domain are not rejected: the polynomial piece of the end span is simply
continued (on the line from $(0,0,0)$ to $(1,0,0)$, \code{c(1.5)} is $(1.5,0,0)$). Do not rely
on this -- use \apiref{extend} (\cref{sec:curve-extend}) for a controlled continuation. At an
interior knot the derivatives are those of the span to the right of the knot.

\subsection{Shape-preserving knot and degree operations}

These operations change the representation of a curve but not its shape (except where a
tolerance is given). They are the classical NURBS algorithms of Piegl \& Tiller.

\begin{description}[style=nextline]
\item[\code{insert_knot(u, times=1)}] Inserts the knot $u$ \code{times} times (A5.1). Each
  insertion adds one control point. The multiplicity cannot exceed the degree.
\item[\code{refine_knots(knots)}] Inserts all values of the list at once (knot refinement, A5.4);
  faster than repeated insertion.
\item[\code{try_remove_knot(u, times=1, tolerance=1e-10)}] Removes the knot $u$ up to \code{times}
  times as long as the curve changes by at most \code{tolerance} (A5.8) and returns the tuple
  \code{(curve, removed)} with the number of removals actually done.
\item[\code{remove_knot(u, times=1, tolerance=1e-10)}] As above, but raises
  \code{GeometryError} if the knot cannot be removed the requested number of times.
\item[\code{elevate_degree(times=1)}] Raises the degree by \code{times} without changing the
  shape (A5.9). The result is clamped.
\item[\code{reduce_degree(tolerance=1e-10)}] Lowers the degree by one if the curve can be
  represented within \code{tolerance}; raises \code{GeometryError} otherwise. The check is
  rigorous: the reduced curve is elevated back and compared with the original. Undoing a degree
  elevation always succeeds.
\item[\code{decompose_bezier()}] The list of Bézier segments (one per non-empty knot span), each
  on its original parameter interval.
\item[\code{clamped()}] The same curve with a clamped knot vector (for curves with an unclamped
  vector, e.g. from periodic data).
\end{description}

Tolerances of knot removal and degree reduction are \emph{Cartesian} deviations in model units.
For rational curves they are checked conservatively on the homogeneous control points, so a
removal may be refused although the true deviation would just be acceptable.

\subsection{Splitting, sub-curves, reversal and transformation}

\begin{description}[style=nextline]
\item[\code{split(u)}] Splits at $u$ and returns the tuple \code{(left, right)}. Both pieces keep
  their parameter intervals $[a, u]$ and $[u, b]$; they are clamped.
\item[\code{sub_curve(u0, u1)}] The piece between two parameters, on the interval $[u_0, u_1]$.
\item[\code{reversed()}] The same point set traversed in the opposite direction: $r(a) = c(b)$.
\item[\code{transformed(transform)}] The curve under an affine \apiref{Transform}
  (\cref{ch:transforms}). Because NURBS are affinely invariant, only the control points are
  transformed; the result is exact.
\end{description}

\apiref{reparametrize}\code{(curve, a, b)} maps the domain linearly to $[a, b]$ (the knots are
scaled and shifted, the shape is unchanged). This is useful when curves with a common
parametrisation are required, e.g. for \code{ruled}.

\subsection{Example}

\example{05_nurbs_curve.py}

Output:

\begin{lstlisting}[style=text]
degree 3 domain (0.0, 1.0) clamped True rational False
C(0.5)  = Vec3(4.170824, 1.964664, 0.000000)
C'(0.5) = Vec3(5.895692, -2.613379, 0.000000)  C''(0.5) = Vec3(-7.709751, 15.532880, 0.000000)
unit tangent Vec3(0.914210, -0.405241, 0.000000)
samples
 [[0.     0.     0.    ]
 [2.3843 2.4248 0.    ]
 [4.1708 1.9647 0.    ]
 [5.5434 2.158  0.    ]
 [8.     2.     0.    ]]
from_points knots [0.0, 0.0, 0.0, 0.0, 0.3333, 0.6667, 1.0, 1.0, 1.0, 1.0]
insert: 8 ctrl, refine: 10 ctrl, max deviation 3.6e-15
removed 2 times -> 6 ctrl
elevated degree 5 n_ctrl 12 -> reduced back to degree 3
reduce_degree refused: reduce_degree: curve cannot be degree-reduced within tolerance 0.00000 ...
Bezier segments: [(0.0, 0.3), (0.3, 0.6), (0.6, 1.0)]
split domains (0.0, 0.5) (0.5, 1.0)
sub_curve domain (0.2, 0.7)
reversed: r(0) == c(1): True
reparametrized domain (0.0, 10.0)
quarter circle: weights [1.0, 0.707107, 1.0] |C(0.3)| = 1.0
\end{lstlisting}

The last lines build a quarter circle by hand: three control points, the middle weight
$\cos 45^\circ = \sqrt2/2$, and every point has distance exactly 1 from the origin.

\section{Joining curves}\label{sec:join}

\apiref{join}\code{(a, b, tolerance=1e-10)} appends curve \code{b} to the end of curve \code{a}
and returns one curve. The end of \code{a} and the start of \code{b} must agree within
\code{tolerance}; a remaining gap is closed at its midpoint. Internally the degrees are unified by
degree elevation, the weights of \code{b} are scaled so that the joint weights agree, the domain
of \code{b} is shifted to start where \code{a} ends, and the joint knot is removed as far as
this is exact. Consequently \code{join(*c.split(u))} reproduces \code{c.clamped()}, and joining a
line with a tangent arc gives a curve that is $G^1$ at the joint.

\begin{warning}
\code{join} does not reverse curves. If \code{b} starts at the \emph{start} of \code{a}, reverse
one of them first. Only touching curves are joined: a gap larger than \code{tolerance} raises
\code{GeometryError}.
\end{warning}

\section{Extending curves}\label{sec:curve-extend}

\begin{lstlisting}[style=text]
extend(curve, end, length, type=Extension.Natural) -> NurbsCurve
\end{lstlisting}

\apiref{extend} lengthens a curve at one end by the given \emph{arc length}; \code{end} is
\code{CurveEnd.Start} or \code{CurveEnd.End}. The original parameters are kept -- the old
domain is a sub-interval of the new one -- so points computed on the original curve remain
valid. Two kinds of extension are available:

\begin{description}[style=nextline]
\item[\code{Extension.Natural}] Continues the polynomial (or rational) piece of the end span. The
  extension is $C^\infty$ at the old end, and conics stay conics: an arc is extended along its
  circle. Extensions that would require non-positive weights are rejected.
\item[\code{Extension.Linear}] Appends a straight segment along the end tangent with the same
  parameter speed as the curve at its end. The result is $C^1$ for non-rational curves and $G^1$
  otherwise.
\end{description}

The same function extends surfaces across one of their sides (\cref{sec:surface-extend}).

\section{Offset curves}\label{sec:offset-curve}

\begin{lstlisting}[style=text]
offset_curve(curve, distance, plane_normal=(0, 0, 1), tolerance=1e-7) -> NurbsCurve
\end{lstlisting}

\apiref{offset_curve} returns the planar offset
$\mathbf o(t) = \mathbf c(t) + d\,\widehat{\mathbf n \times \mathbf c'(t)}$ of a curve lying in a
plane with normal $\mathbf n$. A positive distance offsets to the \emph{left} of the direction of
travel seen from the tip of \code{plane_normal}; for a counter-clockwise circle this is the
inside. The exact offset of a NURBS is not a NURBS (the unit normal contains a square root), so
the offset is a cubic B-spline fitted to the tolerance. The parametrisation is kept: $\mathbf
o(t)$ approximates the offset point of $\mathbf c(t)$.

The function raises \code{GeometryError} where the offset would develop a cusp, i.e.\ where
$|d|$ reaches the radius of curvature on the concave side, where the curve has a corner (a
tangent discontinuity that would open a gap), where $\mathbf c'$ vanishes, and if the curve is not
planar in the given plane. Global self-intersections of the offset (bottlenecks narrower than
$2|d|$) are not detected.

\section{Fillets}\label{sec:curve-fillets}

\subsection{Fillet between two curves}

\begin{lstlisting}[style=text]
fillet_curves(a, b, radius, plane_normal=(0, 0, 1), side_a=0, side_b=0,
              trim=True, corner=None, tolerance=1e-7) -> CurveFillet
\end{lstlisting}

\apiref{fillet_curves} computes the circular fillet of the given radius between two coplanar
curves that intersect. The fillet arc is exact (a rational arc). The sides select which of the
four possible fillets at a crossing is meant: \code{+1} places the centre to the left of the
curve's direction seen from \code{plane_normal}, \code{-1} to the right, and \code{0} chooses
automatically the corner between the \emph{kept parts} -- on each curve, the longer side of the
intersection point. For a polyline corner $a \to X \to b$ this is exactly the inner corner.
\code{corner} selects the intersection point if the curves cross more than once.

The result \apiref{CurveFillet} contains:

\begin{tabularx}{\linewidth}{@{}l>{\raggedright\arraybackslash}X@{}}
\toprule
Field & Meaning \\
\midrule
\code{arc} & the exact fillet arc, running from \code{a} to \code{b} \\
\code{center} & centre of the arc \\
\code{param_a}, \code{param_b} & parameters of the tangent points on \code{a} and \code{b} \\
\code{point_a}, \code{point_b} & the tangent points \\
\code{trimmed_a}, \code{trimmed_b} & the kept parts of the curves, oriented so that
  \code{trimmed_a}, \code{arc}, \code{trimmed_b} form a $G^1$ chain (with \code{trim=True}) \\
\bottomrule
\end{tabularx}

Two lines at an angle $\theta$ give the setback $r/\tan(\theta/2)$ from the corner; in the
example below ($\theta = 90^\circ$, $r = 0.5$) the tangent points are 0.5 from the corner.

\subsection{Rounding all corners}

\apiref{fillet_corners}\code{(curve, radius, tolerance=1e-7)} rounds every corner of a planar
curve -- typically a polyline, polygon or rectangle outline -- with exact arcs and returns a
single $G^1$ curve. A rectangle $4\times2$ with corner radius 0.5 has the length
$12 - 8\cdot0.5 + 2\pi\cdot0.5$.

\subsection{Example}

\example{05_curve_operations.py}

Output:

\begin{lstlisting}[style=text]
join: degree 2, domain (0.0, 2.0), length 3.570796 (2 + pi/2 = 3.570796)
extend natural: domain (0.0, 1.4038182881544163), end Vec3(-0.479426, 0.877583, 0.000000), |end| = 1.000000000000
extend linear : domain (0.0, 1.3535533905932737), end Vec3(-0.500000, 1.000000, 0.000000)
original parameters kept: 0.0e+00
offset +0.2: radius 0.8 within 2.4e-08, degree 3, 137 control points
offset -0.2: length 7.53982237 (2 pi 1.2 = 7.53982237)
offset 1.5 refused: offset_curve: |distance| = 1.500000 reaches the radius of cu ...
fillet: centre Vec3(3.500000, 0.500000, 0.000000), tangent points Vec3(3.500000, 0.000000, 0.000000) / Vec3(4.000000, 0.500000, 0.000000)
        params on a/b: 0.8750 / 0.1667, arc length 0.785398 (pi/4 = 0.785398)
G1 chain length 6.785398 (3.5 + pi/4 + 2.5 = 6.785398)
rounded rectangle: length 11.141593 (12 - 8*0.5 + 2 pi 0.5 = 11.141593)
\end{lstlisting}

Note how the natural extension of the quarter arc stays on the unit circle ($|end| = 1$),
whereas the linear extension leaves it along the tangent at $(0,1,0)$. The offset of the unit
circle by $+0.2$ is a circle of radius 0.8 within $2.4\cdot10^{-8}$, well inside the requested
$10^{-7}$; it needs 137 control points because the cubic fit is done per Bézier segment of the
rational circle.

\section{Curve analysis}\label{sec:curve-analysis}

\subsection{Differential properties}

\begin{description}[style=nextline]
\item[\apiref{curvature}\code{(curve, u)}] $\kappa = |\mathbf C' \times \mathbf C''| / |\mathbf C'|^3$, the
  reciprocal of the radius of curvature.
\item[\apiref{torsion}\code{(curve, u)}] $\tau = (\mathbf C' \times \mathbf C'') \cdot \mathbf
  C''' / |\mathbf C' \times \mathbf C''|^2$; zero for planar curves.
\end{description}

For a helix of radius 1 and pitch 2 the exact values are $\kappa = 1/(1+c^2)$ and $\tau =
c/(1+c^2)$ with $c = 2/(2\pi)$. The fitted helix with tolerance $10^{-9}$ reproduces them to
about $5\cdot10^{-7}$ -- derivatives of an approximation are less accurate than its points.

\subsection{Arc length}

\begin{description}[style=nextline]
\item[\apiref{length}\code{(curve, rel_tol=1e-10)}] the total length,
\item[\apiref{arc_length}\code{(curve, u0, u1, rel_tol=1e-10)}] the length between two parameters,
\item[\apiref{parameter_at_length}\code{(curve, s, rel_tol=1e-10)}] the parameter at which the
  arc length measured from the start of the domain equals $s$.
\end{description}

Lengths are computed by adaptive Gauss--Kronrod quadrature (7/15 points) on every knot span to
the relative accuracy \code{rel_tol}. \code{parameter_at_length} inverts the arc length by a
safeguarded Newton iteration. Use it to place points at equal distances along a curve, e.g.\ for
the positions of features on a circle.

\subsection{Closest points}

\begin{description}[style=nextline]
\item[\apiref{closest_point}\code{(curve, point, tolerance=Tolerance())}] the \emph{global}
  closest point (point inversion): a global candidate search over the whole curve followed by
  Newton refinement.
\item[\apiref{closest_point_local}\code{(curve, point, u0, tolerance=Tolerance())}] a Newton
  iteration from the start value \code{u0}. It is much faster but finds only the local minimum
  near \code{u0}; it stops at the ends of the domain and does not wrap around closed curves.
\end{description}

Both return a \apiref{CurvePoint} with the fields \code{u} (parameter), \code{point} and
\code{distance}. The same functions exist for surfaces (\cref{ch:surfaces}).

\subsection{Knot utilities}

\begin{description}[style=nextline]
\item[\apiref{clamped_uniform_knots}\code{(degree, n_ctrl)}] the clamped uniform knot vector on
  $[0,1]$ for the given degree and number of control points.
\item[\apiref{unique_knots}\code{(knots)}] the distinct knot values and their multiplicities as
  two lists -- the format used by gmsh, OpenCASCADE and STEP.
\item[\apiref{validate_knots}\code{(knots, degree, n_ctrl)}] raises \code{GeometryError} with a
  message if the knot vector is invalid for the given degree and number of control points;
  returns \code{None} otherwise.
\item[\apiref{make_compatible}\code{(curves)}] returns copies of all curves with a common degree
  (the highest), a common knot vector (the union of all knots) and the domain $[0,1]$. Lofts,
  Coons patches and Gordon surfaces do this automatically; call it yourself when you need
  compatible curves, e.g.\ to blend control points.
\end{description}

\subsection{Example}

\example{05_curve_analysis.py}

Output:

\begin{lstlisting}[style=text]
curvature 0.90799984  exact 0.90800033
torsion   0.28902557  exact 0.28902548
length 12.566370614359  (4 pi = 12.566370614359)
arc_length(0, 0.125) = 1.5707963268
parameter at s = pi: u = 0.250000000000, point Vec3(-2.000000, 0.000000, 0.000000)
closest point: u = 0.875000, point Vec3(1.414214, 1.414214, 0.000000), distance 2.455491
expected distance sqrt((3 sqrt2 - 2)^2 + 1) = 2.455491
local search from u0 = 0.8: u = 0.875000, distance 2.455491
local search from u0 = 0.1: u = 0.000000, distance 3.316625
clamped_uniform_knots(3, 7): [0.0, 0.0, 0.0, 0.0, 0.25, 0.5, 0.75, 1.0, 1.0, 1.0, 1.0]
unique_knots: ([0.0, 0.25, 0.5, 0.75, 1.0], [3, 2, 2, 2, 3])
validate_knots: knot vector must be non-decreasing
compatible: degree 2, 5 ctrl, domain (0.0, 1.0)
compatible: degree 2, 5 ctrl, domain (0.0, 1.0)
\end{lstlisting}

The local search started at $u_0 = 0.1$ runs to the start of the domain and stops there with a
distance of 3.32 instead of the true minimum 2.46 -- an illustration of why the global
\code{closest_point} is the safe default.
